\section{确定性采样：$\sigma_t = 0$ 的特殊情况}

\subsection{从随机过程到确定性映射}

当 $\sigma_t = 0$（即 $\eta = 0$）时，DDIM 的采样公式退化为：

$$
\mathbf{x}_{t-1} = \sqrt{\bar{\alpha}_{t-1}}\,\hat{\mathbf{x}}_0(\mathbf{x}_t) + \sqrt{1-\bar{\alpha}_{t-1}}\cdot\frac{\mathbf{x}_t - \sqrt{\bar{\alpha}_t}\,\hat{\mathbf{x}}_0(\mathbf{x}_t)}{\sqrt{1-\bar{\alpha}_t}}
$$

这个更新规则是完全确定性的——没有任何随机噪声注入。这意味着：

\begin{importantbox}{确定性采样的核心性质}
在 DDIM 确定性采样模式下：
\begin{enumerate}
    \item 从\textbf{相同的初始噪声} $\mathbf{x}_T$ 出发，经过整个逆过程后\textbf{一定得到相同的输出} $\mathbf{x}_0$
    \item 噪声空间 $\mathbf{x}_T$ 到数据空间 $\mathbf{x}_0$ 之间建立了一个\textbf{确定性的双射映射}
    \item 生成结果完全可复现：只需记住初始噪声 $\mathbf{x}_T$
\end{enumerate}
这与 DDPM 的随机采样形成鲜明对比——DDPM 中即使从相同的 $\mathbf{x}_T$ 出发，由于每步都注入随机噪声，最终生成的结果每次都不同。
\end{importantbox}

\subsection{确定性采样的实际意义}

确定性采样带来了若干重要的实际好处：

\paragraph{可复现性。} 在实际应用中（如 Stable Diffusion），用户生成了一张满意的图片后，若想再次生成完全相同的结果，在随机采样模式下几乎不可能。而使用 DDIM 确定性采样，只需记录初始噪声种子（seed），即可精确复现任意生成结果。

\paragraph{语义隐空间。} 确定性映射意味着 $\mathbf{x}_T$ 空间中的每个点都唯一对应数据空间中的一张图片。这为隐空间中的操作（如插值、编辑）提供了基础。

\paragraph{与 Neural ODE 的联系。} 确定性 DDIM 可以被理解为 neural ODE 的一种离散化：

\begin{knowledgebox}{DDIM 与概率流 ODE 的关系}
当 $\sigma_t = 0$ 时，DDIM 的迭代更新对应于以下常微分方程（ODE）的 Euler 离散化：
$$
d\mathbf{x} = \left[f(\mathbf{x}, t) - \frac{1}{2}g(t)^2 \nabla_\mathbf{x} \log q_t(\mathbf{x})\right] dt
$$
这就是 Song et al. (2021) 提出的\textbf{概率流 ODE}（probability flow ODE）。ODE 的解轨迹是确定性的，且保持了与 SDE（随机微分方程）相同的边缘分布。使用更高阶的 ODE 求解器可以进一步减少采样步数。
\end{knowledgebox}

\subsection{本章小结}

$\sigma_t=0$ 的确定性 DDIM 建立了噪声空间到数据空间的确定性映射，实现了生成结果的精确可复现性，并为隐空间操作提供了基础。它可以被理解为概率流 ODE 的 Euler 离散化，这一联系为后续使用更高阶 ODE 求解器加速采样提供了理论依据。


